%% ===== Results (~1,700-2,000 words; ONE main figure per subsection) =====
\section{Results\cb{target ~1,900 words}\wc{results}\secstatus{todo}}\label{sec:results}

\bmhead{A shared ontology and knowledge graph unify metabolic-engineering data\cb{~380 words}\wc{r1}\secstatus{todo}}
TorchCell confronts the data-fragmentation barrier by unifying heterogeneous phenotype
data under a single schema and exposing it as deep-learning-ready datasets. Each
experimental record is typed by an experiment ontology grounded in the Biolink Model and
validated at ingestion---for acyclicity, subtype substitutability, Biolink conformance,
and referential integrity---so that records from unrelated studies share one
machine-checkable structure (Fig.~\ref{fig:torchcell}a; ontology data model in
\suppfig{fig:data-model}). Validated records are emitted as typed
subgraphs into a Neo4j knowledge graph hosted at NCSA, and datasets are recovered as
queries over this shared structure: a data-specific query returns one instance with its
typed neighborhood, while a cross-experiment query joins studies through shared entities
such as a common environment or genotype. This turns modular literature and public
databases into interoperable, queryable datasets
\textit{[placeholder: TorchCell currently spans ${\sim}N$ studies and ${\sim}M$ labeled
instances across fitness, gene interaction, expression, and morphology; dataset scale in
\supptab{tab:datasets}, integrated graphs in
\supptab{tab:databases}]}.

Rather than store a genome per strain, TorchCell factorizes data as a shared reference
cell acted on by a perturbation (Fig.~\ref{fig:torchcell}b). Each experiment is a tuple of
a cell graph, an environment, a perturbation, and a phenotype; the minimal reference is a
genome, over which optional gene, regulatory, protein-interaction, and metabolic networks
add structure. The yeast knockout library, for instance, is a single reference with
${\sim}10^{4\text{--}6}$ deletion strains, each represented as a perturbation over that
reference rather than as an independently rebuilt graph. A mechanistic join-and-aggregation
step then merges compatible records---recasting lethality to fitness, or pooling
equivalent genotypes measured under different assays---enlarging datasets while shrinking
the label vocabulary the model must decode (\suppfig{fig:conversion-dedup-agg}).

This organization exposes a favorable information budget (Fig.~\ref{fig:torchcell}c). The
cell's molecular parts---DNA, RNA, protein, and small molecules---are \emph{persistent
entities}: a sequence is the same object in a glucose assay and in an oxidative-stress
assay, and the public archives already hold ${\sim}10^{13}$ compressed bits of them
(\supptab{tab:entity-corpora}). Measured phenotypes are \emph{contingent
observations}, values pinned to one genotype, environment, and assay; the corpus integrated
here amounts to ${\sim}10^{10}$ compressed bits under the same measure. The two therefore
differ by three orders of magnitude in scale and, more consequentially, in reuse: an entity
representation serves every instance in which that entity appears, whereas a phenotype
record constrains the map at one context and no other. We therefore cut the model at the
shared representation, taking entity encoders already trained on the abundant side and
spending the scarce labels only on the perturbation operator and the decoders
(\suppnoteref{note:info-accounting}). Consistent with this budget, classical baselines predict knockout
fitness from almost any representation yet fail on gene interactions (\suppfig{fig:classical-ml})---so fitness alone does not justify a deep model, but the
harder, structured labels do, motivating the architecture that follows.

%% --- Figure 1 (R1): TorchCell database + software overview ---
\begin{figure*}[t]
\tcfig{figures/Fig1-torchcell-overview.pdf}{TorchCell and the Cell Graph Transformer (CGT).
\textbf{Top row --- data.}
(a) Literature and public databases are annotated with TorchCell's Biolink-grounded
experiment ontology and ingested into a Neo4j knowledge graph (hosted at NCSA), queryable
across datasets.
(b) Each experiment is a shared reference cell (genome, gene networks, metabolism, Gene
Ontology) acted on by an experiment-specific perturbation (genotype, environment) with a
measured phenotype; many strains reuse one reference.
(c) Information accounting: molecular parts (DNA, RNA, protein, small molecules) enter as
\emph{persistent entities}---identities that do not change between assays, and whose public
archives carry ${\sim}10^{13}$ compressed bits---and are mapped by encode ($F_\theta$),
perturb ($\mathcal{T}_\psi$), and decode ($\mathcal{R}_\phi$) into \emph{contingent
observations}, phenotypes measured in one genotype, environment, and assay, of which we
integrate ${\sim}10^{10}$ compressed bits. Bit counts are compressed lengths of the
distributed archives---an upper bound on description length, not an information content.
The model is cut at the shared representation $H$. \suppnoteref{note:info-accounting} gives the
full accounting for this panel and states precisely what the bit counts do and do not mean;
\supptab{tab:entity-corpora} reports the measured size of every
archive summed here, with its source and release.
\textbf{Bottom row --- model.}
(d) CGT overview: multimodal cell input $\to$ cell encoder (ENC) $\to$ perturbation
operator (PERT) $\to$ multitask decoder (DEC) $\to$ per-task outputs and loss.
(e) Cell-encoder self-attention block: scaled dot-product attention $\alpha^{(\ell,a)}$
with pre-norm, feed-forward, and residual connections.
(f) Perturbation operator: a soft cross-attention $\beta_{i,t}$ over perturbation tokens
displaces each gene embedding, $h_i^{\mathrm{pert}}=h_i+\Delta h_i$; the same operator
extends from gene deletions to drug and environment perturbations.
(g) Functional equivalence: applying the operator to the reference representation
approximates re-encoding the perturbed graph with a graph network,
$\Phi(G_p)\approx\mathcal{R}_\phi(\mathcal{T}_\psi(H,p))$, so a strain is obtained without
rebuilding or re-encoding its graph (proof in Methods and \suppnoteref{note:equivalence}).
(h) Graph-regularized attention: each biological adjacency $A^{(k)}$ is row-normalized to
$\tilde{A}^{(k)}=(D^{(k)})^{-1}A^{(k)}$ and matched to an attention head by a
KL prior $\Omega_k$, a soft constraint rather than a hard mask (Methods).}{fig:torchcell}
\end{figure*}
\FloatBarrier

\bmhead{The Cell Graph Transformer predicts trigenic gene-gene interactions at state of the art\cb{~380 words}\wc{r2}\secstatus{todo}}
\textit{[FILLER -- R2.]} The Cell Graph Transformer (CGT; architecture in
Fig.~\ref{fig:torchcell}c and Methods) constrains attention with biological interaction
graphs and maps the reference cell through an equivariant perturbation operator. On
trigenic gene-gene interactions (GGI), where classical ML fails (\suppfig{fig:classical-ml}), CGT
reaches Pearson $r = 0.454$ (Spearman $0.421$), outperforming Yeast9 FBA, DCell, and DANGO
(Fig.~\ref{fig:ggi}).

%% --- Figure 2 (R2): Trigenic gene-gene interactions (plan: "Trigenic GI -- can split in two") ---
% Fig 2 is composed in draw.io (source: notes/assets/drawio/Fig2-gene-interactions.drawio)
% and placed true-size with \tcfig. Panel plan, for reference as it is built out:
%   (a) predicted vs actual trigenic interaction; (b) nested cell representation + iBioFoundry
%   DBTL loop; (c) state of the art vs Yeast9 FBA / DCell / DANGO; (d) accuracy vs |tau|;
%   (e) data & model scaling; (f) graph-regularization ablation. Split into two figures if
%   it gets crowded. Update the caption panel-by-panel as each is finalized.
\begin{figure*}[t]
\tcfig{figures/Fig2-gene-interactions.pdf}{TorchCell for gene--gene interaction prediction
and strain design. (a) Higher-order genetic interactions are defined by subtracting
lower-order effects: to isolate a $k$-way interaction you must peel off all interactions
from subcollections of the same genes. For a trigenic interaction, the observed
triple-mutant fitness is compared against the fitness expected from the single- and
double-mutant terms,
$\tau_{ijk}=f_{ijk}-f_if_jf_k-\varepsilon_{ij}f_k-\varepsilon_{ik}f_j-\varepsilon_{jk}f_i$,
with digenic interactions $\varepsilon_{ij}=f_{ij}-f_if_j$.
(b) The nested cell representation---genome, gene networks, metabolism, and
environment---feeds TorchCell and a Learn--Design--Build--Test loop with the NSF
iBioFoundry.
(c) Predicted versus measured trigenic gene--gene interaction.
(d) State of the art on trigenic interactions: Pearson $r$ for Yeast9 FBA ($0.0006$),
DCell ($0.157$), DANGO ($0.367$), and TorchCell ($0.454$).}{fig:ggi}
\end{figure*}
\FloatBarrier

\bmhead{One embedding, many phenotypes: expression, morphology, and fitness\cb{~380 words}\wc{r3}\secstatus{todo}}
\textit{[FILLER -- R3.]} The same embedding predicts knockout expression ($r = 0.543$),
single-knockout morphology ($r = 0.619$), and fitness (Fig.~\ref{fig:multitask};
expression diagnostics in \suppfig{fig:expr-diag}).

%% --- Figure 3 (R3): expression, morphology, fitness ---
% Fig 3 composed in draw.io (source:
% notes/assets/drawio/Fig3-systems-biology-fitness-expression-morphology.drawio) and placed
% true-size with \tcfig. Update the caption panel-by-panel as each is finalized.
\begin{figure*}[t]
\tcfig{figures/Fig3-systems-biology-fitness-expression-morphology.pdf}{One latent cell
embedding generalizes across systems-biology phenotypes---expression, morphology, and
fitness. (a)--(e).}{fig:multitask}
\end{figure*}
\FloatBarrier

\bmhead{Natural genetic variation versus model-designed perturbations\cb{~380 words}\wc{r4}\secstatus{todo}}
\textit{[FILLER -- R4.]} We compare phenotypes from natural strain genetic variation against
deep-learning-designed gene perturbations across genome and environment
(Fig.~\ref{fig:variation}).

%% --- Figure 4 (R4): natural variation vs DL design ---
% Fig 4 composed in draw.io (source: notes/assets/drawio/Fig4-natural-variation.drawio) and
% placed true-size with \tcfig. Update the caption panel-by-panel as each is finalized.
\begin{figure*}[t]
\tcfig{figures/Fig4-natural-variation.pdf}{Natural genetic variation versus model-designed
perturbations. (a)--(c).}{fig:variation}
\end{figure*}
\FloatBarrier

\bmhead{Drug exposure and robustness of the learned cell representation\cb{~380 words}\wc{r-robust}\secstatus{todo}}
\textit{[FILLER -- R-robust.]} The perturbation operator extends from gene deletions to drug and
environment exposure, and the learned cell representation is robust: predictions stay stable
across dataset splits, tolerate embedding and interaction-graph ablations, and degrade
gracefully as supervision is reduced (Fig.~\ref{fig:robustness}).

%% --- Figure 5 (R-robust): drug exposure + robustness ---
% Fig 5 composed in draw.io (source:
% notes/assets/drawio/Fig5-drug-exposure-and-robustness.drawio.svg) and placed true-size with
% \tcfig. Update the caption panel-by-panel as each is finalized.
\begin{figure*}[t]
\tcfig{figures/Fig5-drug-exposure-and-robustness.pdf}{Drug exposure and robustness of the
learned cell representation. (a)--(f).}{fig:robustness}
\end{figure*}
\FloatBarrier

\bmhead{Extending to metabolism and strain design\cb{~380 words}\wc{r5}\secstatus{todo}}
\textit{[FILLER -- R5.]} Grounding the cell representation in metabolism, CGT recommends gene
targets for production phenotypes and pairs with the UIUC iBioFoundry for iterative design
(Fig.~\ref{fig:metabolism}; inference at scale in \suppfig{fig:inference-scale}).

%% --- Figure 6 (R5): metabolism and metabolic engineering ---
% Fig 6 composed in draw.io (source:
% notes/assets/drawio/Fig6-metabolism-and-metabolic-engineering.drawio.svg) and placed
% true-size with \tcfig. Update the caption panel-by-panel as each is finalized.
\begin{figure*}[t]
\tcfig{figures/Fig6-metabolism-and-metabolic-engineering.pdf}{Extending TorchCell to metabolism
and metabolic-engineering strain design. (a)--(f).}{fig:metabolism}
\end{figure*}
\FloatBarrier
